\begin{textsample}{NEG Dim 6 – pi – Score 1.00 – pi/pi\_em\_57.txt}  \label{ex:f6_neg_097}
3.4 Sociologia A sociologia é uma ciência que agrega parcela significativa das Ciências Humanas na formação dos jovens, visto que todas as sociedades estão vinculadas às mudanças culturais e sociais, definindo transformações que impactam na vida dos indivíduos, em especial a dos jovens.
Dessa forma, a Sociologia contribui para moldar a personalidade, interesses, projetos de vida e relações com o mundo em que os estudantes vivem.
Como componente curricular, a Sociologia surgiu a partir da década de 1920, e enfrentou percalços até se consolidar nas instituições escolares.
Em 2009, tornou -se componente obrigatório nos currículos do ensino médio, estando regulamentada pelo Conselho Estadual de Educação ( CEE ) do Piauí, através da Resolução No CEE/PI no 111/2009, que fixou o ensino de Sociologia no Ensino Médio no âmbito do Sistema de ensino do Estado.
Nesse novo contexto de currículo o componente de Sociologia se articula com os conhecimentos alinhados às competências gerais da BNCC e às competências específicas da área de Ciências Humanas e Sociais Aplicadas.
A Sociologia torna -se essencial para o desenvolvimento das habilidades necessárias para articular a construção do conhecimento e saberes necessários no Século XXI.
Visando garantir os direitos dos estudantes do Piauí a fim de desenvolverem competências e habilidades esperadas na educação básica, o currículo traz na sua estrutura de componentes a Sociologia como meio dos estudantes interagirem com o meio social, pois a partir dela os indivíduos são capazes de refletirem acerca dos desafios da atualidade e proporcionar uma reflexão sobre experiências e conhecimentos inerente a sociedade.
O ensino de Sociologia deve atender preferencialmente os vínculos comunitários, organizações, e movimentos com base nas ideias propagadas pela BNCC de justiça, solidariedade, autonomia, equidade, liberdade de pensamento e de escolha, compreensão e o reconhecimento das diferenças e da interculturalidade.
Nessa perspectiva, seus conteúdos devem contemplar uma abordagem inter, trans e multidisciplinar com vistas a contemplar objetos do conhecimento específicos das demais áreas, temáticas de conhecimento geral e específico visando ao domínio, ao engajamento profissional e ao projeto de vida.
Nesse sentido, a Sociologia visa incentivar os sujeitos a participarem do debate público de forma crítica, respeitando diferentes posições e fazendo escolhas alinhadas ao exercício profissional, com liberdade, autonomia, consciência crítica e responsabilidade social, considerando o que preconiza a BNCC de entrelaçamento com o projeto de vida envolta de questões partindo do local até o global.
O objetivo da Sociologia no ensino médio é oferecer para o estudante múltiplos olhares sobre sua construção como indivíduo social, bem como acerca das múltiplas estruturas que sustentam as sociedades.
Para isso, a sociologia se baseia em dois princípios básicos do pensar : 1 ) a desnaturalização, que mostra que às explicações do senso comum tendem a naturalizar os fatos sociais – por exemplo, “ é da natureza da sociedade ”, “ sempre se deu assim ”, “ é da vontade de Deus ” – argumentos que o trabalho científico desconstrói a partir da observação das relações sociais e processos históricos ) e 2 ) o estranhamento, que corresponde ao princípio elementar de todas a ciência : os fatos não se dão a conhecer imediatamente, mas por meio de questionamentos e problematizações.
Dessa forma, a discussão de temáticas que perpassam a área como cultura, socialização, trabalho, juventudes, Estado, cidadania, educação, movimentos sociais, classes sociais, violência, identidades, desigualdades sociais, castas, fato social, dominação, alteridade dentre outros, tornam -se fundamental para aproximar os estudantes do campo da política, da cultura, da ciência do Estado e da sociedade, em âmbito local, regional, nacional e mundial, visando despertar o pensamento crítico acerca da realidade.
A Sociologia, em diálogo com os demais componentes da área Ciências Humanas e Sociais Aplicadas, através das práticas de reflexão sobre os discursos do senso comum em diálogo com os saberes científicos e as realidades investigadas, fomenta entre os estudantes práticas de protagonismo individual e social.
Para protagonizar, os indivíduos precisam aprender por meio de experiências pedagógicas sistemáticas, coerentes, críticas, dialógicas, criativas e criadoras.
Nessa perspectiva, o currículo do Estado do Piauí tenta aproximar o estudante de forma divertida, criativa e comprometida com a educação mediante debate e leituras que impulsionam o processo de ensino aprendizagem, considerando os jovens como atores políticos na educação.
Por fim, a sociologia é uma disciplina que estuda a sociedade e sua organização, como os indivíduos se ligam aos grupos e como interagem uns com os outros.
A sociedade contemporânea exige cada vez mais da escola a oferta do desenvolvimento de competências e habilidades entre os estudantes, para que estes tenham condições de agir com consciência e responsabilidade enquanto cidadãos.
A sociologia torna -se fundamental, para a construção deste comportamento de uns com os outros, sendo um ponto forte para esses jovens aprenderem a conviver e a trabalhar em coletividade.

% matched lemmas: 
\end{textsample}
